\section{Introduction}
\label{sec:intro}

Enterprises accumulate process regulations, failure analyses, and operational
lessons as unstructured documents spread across shared drives. Answering a
concrete engineering question --- \emph{why did edge delamination appear in an
EL inspection of a laminated photovoltaic module?} --- requires locating the
one report among thousands that records the decision, reading it, and citing
it. Retrieval-augmented generation (RAG) addresses the location step by
embedding documents and the query into a shared vector space
\cite{rag-lewis,dpr}, and modern agentic stacks let a large language model
(LLM) orchestrate retrieval tools iteratively \cite{mcp,singh-agentic-rag}.
Yet the dominant ranking signal --- dense cosine similarity --- is a measure
of \emph{topic relatedness}, not of \emph{evidence presence}. A document that
merely discusses similar materials can outrank the one document that contains
the answer, and a fluent generator can readily synthesize an answer from the
wrong source. The result is the well-documented hallucination and
misgrounding behavior of LLM pipelines \cite{siren}.

This paper argues that a primitive most deployed RAG stacks leave implicit
is an \emph{explicit content-verification gate}: before a candidate document
is allowed to support an answer, its actual text must be fetched and judged
against the query by a model whose only job is to answer ``does this text
contain direct evidence for this question?'' Vectors are fast; only the
content can confirm. We built QDCVR (Query-Driven Content-Verified Retrieval), a
self-hosted platform that makes this primitive the backbone of a complete
knowledge-base system --- content-routed organization at ingestion, a
97-tool MCP server as the machine interface \cite{mcp}, and three retrieval
lanes that all terminate at the same verification gate --- and, because the
primitive is itself a model, subjects it to the same scrutiny as any model
(Section~\ref{sec:inflation}).

The three lanes embody a deliberate latency/recall trade-off. The
\emph{Search} lane is a vector fast lane: rewrite the query, retrieve a wide
candidate net, and submit citation-only references to the verification gate. The
\emph{Librarian} lane is a hierarchical navigator that walks
library$\to$document catalogs, routes enumeration-style queries to a machine
scanner, judges survivors in small batches, and reads only what survived ---
all within a hard budget of twelve tool calls. The \emph{Hybrid} lane
collapses vector and catalog channels into a single server-side orchestrated
call that merges candidates, judges them, applies a relative cut with a
top-5 floor, and returns full text for the top documents. All lanes fail
closed: zero verified evidence yields an explicit not-found report, never a
speculative answer.

We assess the platform with two experiment tiers on a deployed industrial
corpus of five knowledge bases spanning photovoltaic materials, module
failure analysis, plant O\&M, and discrete-manufacturing telemetry
(42 documents): a tool-layer
comparison of seven retrieval pipelines plus two derived configurations
on sixteen gold-annotated questions (plus a sixteen-question paraphrase
grid), and
16 end-to-end agentic runs that answer four questions through the platform's
public chat API in four arms (three lanes plus a bare agent without
knowledge-base tools), followed by a robustness tier: leave-one-out margin
generalization, a second verification signal replayed over the logged
candidate pools, and graded abstention probes. The contributions are:

\begin{enumerate}
\item \textbf{A verifiable-retrieval architecture} (Sections
\ref{sec:system}--\ref{sec:lanes}): content-verified retrieval as a
platform primitive whose measured value lies in auditability, per-citation
provenance, and fail-closed semantics --- with three interchangeable lanes
under enforced budgets and an engine-agnostic harness layer that exposes
everything as MCP tools.
\item \textbf{A two-tier empirical analysis} (Sections
\ref{sec:setup}--\ref{sec:results}): at the tool layer, BM25 and reciprocal-rank
fusion are strong baselines on a curated same-domain corpus, and dense
retrieval misses exactly one question (Q16) through a near-duplicate
confuser family; the
verification gate's re-ranking \emph{reduces} top-5 hit rate from 0.94 to
0.44 when confusers
--- near-miss same-domain distractors --- dominate the candidate pool. At
the agent layer, all knowledge-base-grounded lanes ground their answers in
gold documents while end-to-end latency is almost entirely LLM turns
(machine retrieval is sub-second).
\item \textbf{A quantified verifier-failure analysis} (Section
\ref{sec:inflation}): an operational quantification, inside a deployed
verified-retrieval pipeline, of LLM-judge miscalibration --- whose
general causes (score compression, self-preference) are documented in the
judging literature \cite{zheng-llmjudge,panickssery-selfpref} but whose
pipeline-level consequences we decompose and measure: same-domain evidence
inflation (non-gold median 0.86 versus gold 0.88; Mann--Whitney $p = 0.26$,
i.e.\ no separation), a verdict role asymmetry (verdicts are safe as a
binary filter --- Hit@5 restored to 0.94, sign test $p = 0.008$ --- and
unsafe as an ordering), a threshold that never abstains out of corpus,
and robustness of both the failure and the repair to question paraphrase
and to the choice of verification signal.
Mitigations (hierarchical reading, relative cuts with floors, calibration)
are scoped to what this decomposition supports.
\item \textbf{Operational lessons from production use} (Section
\ref{sec:discussion}): three documented incidents, each with its detection
symptom and fix --- resident judge workers (batch judging 10--12\,s
$\to$ 1--2\,s), vector-index segment corruption surfacing as honest
``success with zero results'' failures, and GPU memory hygiene between
embedding and judging workloads.
\end{enumerate}

A preliminary demo describing the platform's organization and retrieval
layers was accepted as a CIKM~2026 demo \cite{qdcvr-demo}. This paper extends
it with three components that do not appear in the demo: (i)~the multi-lane
agentic retrieval methodology with enforced per-lane tool budgets; (ii)~the
verifier-failure analysis --- quantified same-domain inflation, the
filter/re-ranker role decomposition, and its configuration-level repair;
and (iii)~the two-tier empirical evaluation on the deployed corpus with all
per-run logs released.

\subsection{Research questions and reading guide}
\label{sec:rq}
The paper is organized as a design-to-evidence arc built around four research
questions, each answered by one evidence tier and reported in one place
(Table~\ref{tab:ataglance}):

\begin{description}
\item[RQ1 (Architecture).] Can content verification be \emph{enforced} as a
platform-wide primitive --- that is, can a deployed knowledge-base platform
guarantee that no citation enters an answer without its content having been
judged, across heterogeneous retrieval strategies and independent agent
harnesses?
\item[RQ2 (Ranking trade-off).] What does the verification gate \emph{cost}
and what does it \emph{buy} at the ranking layer --- and does the answer depend
on the role the verdicts play (a filter, an ordering, or a threshold)?
\item[RQ3 (Verifier generality).] Is the verifier's failure mode an artifact
of one checkpoint, one domain, one language, or one corpus size, or is it
structural?
\item[RQ4 (Operational contract).] Where does honest abstention actually live,
and where does end-to-end latency go?
\end{description}

\begin{table}[!htbp]
\centering
\footnotesize
\setlength{\tabcolsep}{3pt}
\caption{This paper at a glance: each research question, the evidence tier
that answers it, and the headline finding.}
\label{tab:ataglance}
\begin{tabular}{p{0.75cm}p{4.0cm}p{2.6cm}p{5.3cm}}
\toprule
RQ & Question & Evidence tier & Headline finding \\
\midrule
RQ1 & Is verification an enforceable platform primitive? & Architecture and
lanes (\S\ref{sec:system}--\S\ref{sec:impl}) & Yes --- server-side tool
whitelists and fail-closed contracts make it a property of the enforcement
layer, not of prompt discipline \\
RQ2 & What does the gate cost and buy? & Tool layer, nine configurations
(\S\ref{sec:retrieval}--\S\ref{sec:inflation}) & Costs $<$1\,s; as a
\emph{re-ranker} it loses Hit@5 ($0.94\!\to\!0.44$), as a \emph{filter} it is
neutral ($0.94$; sign test $p=0.008$) --- its measured value is per-citation
provenance, not ranking quality \\
RQ3 & Is the verifier failure structural? & Second signal, span-NLI, symbolic
check, second corpus (\S\ref{sec:inflation}, \S\ref{sec:robustness}) &
Structural: the loss tracks confuser density, not checkpoint, domain,
language, or scale \\
RQ4 & Where do honesty and latency live? & Out-of-corpus and graded probes,
turn accounting (\S\ref{sec:abstention}--\S\ref{sec:composition}) & Honesty
lives in the reading--answer layer, never in the threshold; latency is
LLM-turn-dominated \\
\bottomrule
\end{tabular}
\end{table}

The remainder of the paper follows the arc.
Sections~\ref{sec:system} and \ref{sec:lanes} present the architecture and
the three retrieval lanes and answer RQ1 (Section~\ref{sec:trace} first walks
one real question end to end so the components are concrete before the
measurements begin); Section~\ref{sec:impl} reports the engineering that makes
the contracts enforceable; Sections~\ref{sec:setup}--\ref{sec:results} carry
the two experiment tiers plus the robustness tier and answer RQ2--RQ4;
Section~\ref{sec:discussion} distills the operational lessons, threats, and
roadmap; Section~\ref{sec:conclusion} concludes.
Table~\ref{tab:modules} in Section~\ref{sec:system} maps every module to its
responsibility; Table~\ref{tab:ataglance} maps every research question to its
evidence.

